\section{Primer on Out-of-Order Processors}
\label{apendixooo}


The out-of-order design we consider is split between a front-end (responsible for speculatively fetching new instructions) and a back-end (responsible for executing out-of-order those instructions, emitting the corresponding memory operations, and correcting wrong predictions to redirect the front-end).
An overview of the Riscy-OOO pipeline is provided in Figure \ref{fig:riscypipeline}.
The front-end is irrelevant to this paper. 
The backend starts when an instruction is enqueued both into the reorder buffer (see below) and the correct ALU/FPU/Memory execution pipeline, via each pipeline's reservation station (see below).
Morally, the reorder buffer tracks when an instruction is ready to commit, while the reservation stations track when an instruction is ready to execute.
Together, the reorder buffer and the reservation stations ensure that instructions are executed when their operands become available, as long as it does not contradict the program order. They are also responsible for killing mispredicted instructions and committing right-path instructions.
\\

\noindent\textbf{Reorder Buffer}
Every cycle, the different execution pipelines notify the ROB when an instruction has finished executing, and the instructions' statuses are updated in the ROB when done with execution.
To guarantee correctness, the effect of an instruction is permanently committed only once it reaches the head of the ROB (guaranteeing that all previous instructions in the program have committed).
When an instruction is required to \textit{execute non-speculatively}, the scheduler waits until that instruction reaches the head of the ROB before notifying the reservation station (see below) that this instruction, which is now the oldest in-flight instruction, is ready for execution. 
\\

\noindent\textbf{Reservation Station}
There are multiple reservation stations for different kinds of instruction: ALU, FPU, and the memory instructions.
At each cycle, the entries in the reservation stations are updated to reflect if their operands are ready (that is, if the dependent register values have been computed by previous instructions).
Following this, the scheduler selects, for each reservation station, the oldest instruction that is ready to execute (has all of its required operands and is not), and issues it for execution to its corresponding execution pipeline.

\noindent\textbf{Overview}
\begin{figure}[t]
    \centering
    \includegraphics[width=\columnwidth]{figures/riscypipeline.png}
    \caption{Overview of the Riscy-OOO pipeline.}
    \label{fig:riscypipeline}
\end{figure}

\section{Hardware Implementation Details}
\subsection{Safe Mode}
\label{app:safemode}
In the baseline out-of-order processor we are hardening, memory instructions are decoded, renamed, and then simultaneously queued in the reorder buffer (ROB) and the memory reservation station (\MemRS{}) (see Appendix~\ref{apendixooo}).
When an instruction's operands are ready, it is dispatched to the memory execution pipeline (\MemExePipeline{}).
Safe mode identifies instructions that access untrusted memory early in the \MemExePipeline{} and delay the translation and processing of the memory accesses until the instruction is no longer speculative, i.e.,  when it reaches the head of the ROB.
These changes require modification of \MemRS{} and \MemExePipeline{}. 
\\

\noindent\textbf{Modifying the Memory Execution Pipeline}
Upon entry into \MemExePipeline{}, checks are performed on memory instructions to decide if the memory operation is safe to translate and then issue or not.
It is important to perform these checks very early as
the next steps in the \MemExePipeline{} touch shared microarchitecture (through address translation and cache hierarchy, for instance) potentially transmitting information to an attacker.
A memory access is deemed safe if it accesses trusted memory OR if it is non-speculative.
Because the enclave trusted memory range is defined in virtual memory, this check can be performed as soon as operands (i.e. virtual address) are available.
If an instruction tries to access untrusted or shared memory, we check that it is non-speculative by conservatively checking if it has reached the head of the ROB.
If the instruction is safe, it is dispatched to the next stage of the \MemExePipeline{}.
If it is not safe, the instruction is squashed and a signal is sent to the \MemRS{} indicating that it should re-dispatch it once it reaches the head of the ROB.
\\

\noindent\textbf{Decorrelating Request Dispatching and Request Dequeuing}
Safe mode requires the \MemRS{} to retain the accesses to shared memory in order to re-dispatch them once they become non-speculative.
The original design performs a \MemRS{} data dispatch and dequeue simultaneously once the request enters the \MemExePipeline{}. 
These two operations now need to occur at different stages in the pipeline: data is first dispatched from the \MemRS{} and is dequeued in the translation execution stage if and only if the request is determined to be secure. 
We modify the interface between \MemRS{} and \MemExePipeline{} to provide the correct interfaces and logic.
This includes augmenting the request data with the corresponding \MemRS{} entry index.
\\

\noindent\textbf{Modifying the Memory Reservation Station}
Our mechanism requires \MemRS{} being able to re-dispatch some instructions once they become non-speculative: untrusted memory loads will indeed typically need to get dispatched twice.
First, several new status bits are added for each \MemRS{} entry, specifically to track if a) a given request is still to be dispatched for the first time, and b) if a given request needs to be re-dispatched. 
The first bit (\textit{to\_dispatch}) is necessary since the entries can now still reside in the \MemRS{} in the time between being dispatched and dequeued (or signaled for redispatch) and should not be selected for dispatch again.
The second bit (\textit{to\_redispatch}) is set by \MemExePipeline{} when it squashes a speculative insecure request. 
We additionally add logic to compare each \MemRS{} entry to the entry at the head of the ROB, and raise the entry's \textit{dispatch\_ready} signal on a match.
Finally, the dispatch logic is modified to use the new status bits. Rather than simply checking for the oldest valid entry whose arguments are ready, an additional condition is added to check if either \textit{to\_dispatch} is set for the entry or \textit{to\_redispatch} and \textit{redispatch\_ready} are set for the given entry.
\\

\noindent\textbf{Disabling the Mechanism for the OS}
Blocking speculation on memory access can have a significant performance impact.
We want to make sure that the OS is still allowed to speculate on its own memory if it chooses not to use our trusted memory mechanism.
As a result, we disable our mechanism when the enclave virtual address range is set to a specific empty range.
Note that the OS could always be modified to take advantage of our trusted memory mechanism and protect itself from speculative side-channel attacks.

\subsection{Reconfigurable LLC Partitioning}
\label{app:dynamiccache}
Static partitioning (as proposed in MI6) is impractical as each of the 64 DRAM regions would be allocated a 16KB cache slice in our 1MB LLC.
This is even smaller than the 64KB necessary for the inclusion of the cores' private L1 caches (D1 and I1), which is dramatic for performance.
Using way partitioning \cite{kiriansky2018dawg} is not possible as we need to instantiate 64 regions in the cache, but the LLC has only 16 ways, and Ironhide's \cite{ironhide} strategy assumes limited reconfiguration.
\\

\noindent\textbf{Hardware-Software Co-Design}
The key is to rely on the (trusted) SM to do most of the heavy lifting in terms of bookkeeping and checking for correct configuration and security.
At the hardware level, we can now design a new, simple, and practical reconfigurable LLC partitioning scheme that only requires minimal hardware changes.
\\

\noindent\textbf{Configuration Protocol}
We build a mechanism for the OS to reconfigure the LLC partitions to improve enclave and other program's performance.
This costly operation is only required when changing the LLC layout, which happens only before launching an application or an enclave with an important memory footprint.
The OS requests a reconfiguration through an SM call while no enclaves are running.
The SM will first check that the new partitioning is valid, doesn't affect regions assigned to the SM.
It then waits for other cores to flush their L1 caches and to enter a blocked state in the SM, making sure no requests will be sent to the modified regions during reconfiguration (which could put the cache in an incoherent state).
The SM will then interact with the hardware to flush the LLC slices that are modified and set the LLC in its new configuration.
\\

\noindent\textbf{Microarchitectural Implementation}
\begin{figure}[t]
    \centering
    \includegraphics[width=\columnwidth]{figures/partitioning.png}
    \caption{LLC set indexing scheme. The function takes an address as an input and output a set index in the LLC slice allocated to its memory region.}
    \label{fig:dynamicpartitioning}
\end{figure}
The configuration for each LLC slice (base and bound of the set-index range) is stored in a big register (1216 bits) called the set-index range table.
The SM can modify the entries of the set-index range table by sending MMIO requests to the cache, which is seen as a new device.
When the LLC receives a request, it will compute the new set index for the corresponding address.
We represent the transformation used in Figure \ref{fig:dynamicpartitioning}.
First, we extract a region ID from the address upper bits and obtain the corresponding index range base and size by looking up the set-index range table \circled{1}.
We then extract the original set index of the address, located in the middle of the address, and tweak it.
We compute the original set-index value modulo the size of the index range \circled{2} and add the base index \circled{3}.
This effectively restricts the new set index to the corresponding index range, partitioning the LLC \circled{4}.
Note that we must extend the length of the tags used by the LLC to include the entire original index in case an index range is of size 1 and all addresses from the same region map to the same set.
\\

\subsection{LLC fine-grain flush}
\label{app:llcflush}
To go hand-in-hand with our reconfigurable LLC partitioning, we also want to be able to selectively flush parts of the LLC.
We create a \emph{Zero Device} in the upper physical address space. 
It is implemented as a simple piece of logic connected to the cache hierarchy at the same level as the DRAM, but that always sends back zeros.
The trick is that the address space covered by the zero device has the same size as the DRAM and is located at addresses that will collide with DRAM addresses in the cache.
Once the SM has placed the machine in a correct waiting state (where all the cores have flushed their L1 and are waiting in the SM), it only needs to access an eviction set in the zero-device address space to flush the cache slice.
This means that flushing a small LLC slice is much faster than flushing a big one.
We look at the performance of our flushing mechanism in the evaluation section (see \ref{sec:evalllcpartitioning}).

\subsection{Relaxed Flushing Policies}
\label{app:relaxedflushing}
The enclave and some parts of the mini-SM do not have secrets to keep from each other.
Specifically, the mini-SM never manipulates sensitive data.
As a result, the mini-SM does not need to flush microarchitectural state when context switching to enclaves.
This reduces the cost of mini-SM calls and opens the door to the use of frequent software or timer interrupts.
To disable flushing selectively, we modify the hardware to deactivate systematic flushing and exposed the SM a register to trigger the flush.
We write the SM to trigger the flushing mechanism on required context switches, but not on most traps and SM-calls from the enclave.
\\

\noindent\textbf{Evaluation}
\label{app:evalrealxedflush}
\begin{figure}[t]
\centering
\includegraphics[width=\columnwidth]{figures/timer_interrupt_coremark.png}
\caption{Normalized Execution Time for Coremark in the presence of periodic timer interrupts. MI6 flushed on every interrupt; Citadel does not.}
\label{fig:coremark_timer}
\end{figure}
To evaluate the performance of the relaxed flushing policy, we run Coremark while triggering periodic timer interrupts at different time intervals to model frequent traps.
We do not perform any computation inside the trap handler and only measure the overhead due to context switching.
The results are provided in Figure \ref{fig:coremark_timer}.
By testing a range of interrupt periods, we confirm that flushing on each trap like in MI6 results in significant overheads.
However, the cost is greatly reduced when using the relaxed flushing policy, with a maximum of 7\% overhead measured for interrupts every 0.4 ms.

\subsection{Miscellaneous}
We statically partition MSHRs between two cores and place a fair round-robin arbiter at the entry of the LLC.
Because these two mechanisms only marginally affect the latency of the LLC for a 2-core design, we observe negligible ($\approx$ 0.1\%) performance impact.

We also expose controls for speculation and predictors to the software through Control Status Registers (CSRs).
Specifically, a program is able to deactivate speculation all together (every instruction waits to reach the head of the ROB), delay the issuing of shared memory accesses to the head of the ROB, deactivate the training of predictors (branch predictors and return address stack) and/or their use (and then it falls back to the predictor $pc + 4$, or straightline speculation).

\section{Software Infrastructure}
\label{app:softwareinfrastructure}

\subsection{Security Monitor}
The SM does not allocate any resources, this task is left to the OS, but it will interpose between each security-sensitive operation (like a resource allocation) and check that the operating system performed the task without violating the high-level security policies (e.g., it is allocating a memory region to an existing enclave).
Its API is called by the OS or an enclave to request ownership of new memory regions, create enclaves, or perform any other operations where it is needed.
As a result, the SM is mostly responsible for bookkeeping operations, performing security checks, and setting up the hardware mechanisms to enforce isolation between the running security domains.
It will not only systematically intervene on any security-sensitive operations, but will also play the role of a small kernel and build the bridge between the OS and the hardware on some non-security-sensitive operations like sending inter-process interrupts (IPI).
The OS and the enclaves interact with the SM through ecalls.
A more detailed description of the SM data structures and API, can be found in \cite{lebedev2019sanctorum}.

\subsection{Boot Sequence}
\noindent\textbf{Secure Bootloader}
To boot an OS such as Linux on top of our SM, several steps are necessary.
First, our Secure Bootloader will be in charge of the first booting stage.
While other cores are waiting,
core 0 will load, measure (hash) the SM binary, and use the SM measurement to generate the SM key pair. 
Once this is done, the Bootloader will use its private keys to sign the SM measurement and endorse the SM public key, creating the first certificate used for secure boot or remote attestation.
Finally, it will copy the SM key pair to SM memory along with the endorsing certificate.
For more details on secure boot, see~\cite{lebedev2018secure}.
It will then wake up the other cores using interprocess interrupts and have them all jump to the SM entry point.
\\

\noindent\textbf{SM Boot}
Now that the SM has been attested and is running, it can start initializing the hardware, set up the data structures required to later perform its duties, and set up the machine is in such a state so that Linux can properly boot in the next stage.
The SM proceeds as follows: each core will clean up its register file and set up a stack.
They will all initialize some of their special registers.
In particular, they will immediately set up a trap handler to be able to catch exceptions and enable software interrupts.
From there, all cores go to sleep except core 0.
Core 0 will set up all private SM data structures that keep track of the ownership of resources.
It will also walk the device tree and set up kernel-related mechanisms such as the clock and interprocess interrupts.
The other cores will then be awakened.
The SM will set up memory protection through the specific registers.
In particular, some memory checks are performed during address translation.
However, Linux uses physical addresses to boot.
To solve this issue, the SM sets up an identity page table, so the virtual address space exposed to Linux is a direct mapping to the physical address space.
Each OS memory access will then go through address translation and memory security checks.
Finally, the SM will have all the cores jumping to the payload entry point at a similar time; in our case, the payload is the Linux image.
\\

\noindent\textbf{Linux Booting and SM Kernel Module} 
We boot Linux using the Spinwait method \cite{dabbelt2017all}.
Once Linux boots, an SM Kernel Module will be loaded.
The SM kernel module is registered as a device in Linux that is accessible to applications through a pseudo-file \texttt{/dev/security\_monitor}.

\subsection{Running an Enclave}

\subsubsection{Key Software Elements}
\noindent\textbf{The SM Kernel Module}
The main SM should not possibly be made unavailable to the OS by a user-level program.
Furthermore, no resource should be allocated to a user-level program without the OS approval.
As a result, to create an enclave or perform other operations,
a user application needs to interact with the SM through the SM Kernel Module.
Because the SM Kernel Module is part of the OS, it can then refuse requests if needed.
\\

\noindent\textbf{Mini-SM}
Because the SM is a shared resource, its might leak sensitive information. For instance, an enclave might call the SM at particular times and in a pattern that might leak sensitive information regarding the enclave's secrets.
As a result, each enclave uses its own private copy of the SM.
This "mini-SM" is smaller than the one accessible to the OS and only contains the code corresponding to the subset of the SM API accessible by the enclave.
It also uses its own private stack.
This mini-SM is located in private enclave memory (not accessible nor visible by the OS) and protected using physical memory isolation (hence it is neither modifiable nor accessible by the enclave).
The enclave interacts with its mini-SM through the ecall API.
Note that some of these ecalls might touch on shared global SM state and locks.
As a result, an enclave should be careful when interacting with the SM and be aware that it might still leak timing information to other programs.
We emphasize that an enclave binary does not go through the SM Kernel Module and directly interacts with its mini-SM through ecalls.
\\

\noindent\textbf{User-Level Application}
The user-level application runs on top of the OS.
It does not benefit from any integrity or privacy guarantees, unlike its enclave counterpart.
On the other hand, it has access to a rich interface with the outside world (like the file system or the network) through all the usual OS services implemented through system calls.
Even if not required for an enclave to work, this makes the user-level application a good ally for an enclave program.
It can, for instance, fetch encrypted private data and forward it to the enclave or perform many other non-security-sensitive tasks.

\subsubsection{Launching An Enclave}
\label{app:launchingprocess}
\noindent\textbf{Allocating Memory}
The user-level application first loads the enclave binary from the file system to user memory.
It will also ask the OS to map the virtual addresses used by the enclave for shared memory to physical memory.
The application should also make sure not to allocate any shared objects within the enclave's virtual memory range.
The application then opens the kernel module device and sends a request through an \texttt{iotcl} system call to request the SM Kernel Module to launch the enclave binary.
\\

\noindent\textbf{Setting Up the Enclave}
Once called, the SM Kernel Module launches a series of ecalls to the SM to setup the enclave.
First, it moves the enclave binary from user memory to kernel memory.
It then allocates two physical memory regions and performs several SM calls so the SM can use one of the regions to store metadata and the second region can be given to the enclave.
After the SM has set up the metadata region, the enclave's metadata can be created and populated.
During the enclave setup process, the SM makes sure to add several arguments used to set up the enclave to the enclave's measurement (see Section \ref{app:attestation} for more details).
The kernel module then asks the SM to copy a version of the mini-SM in enclave memory, initializes the enclave's page table, and loads the enclave's binary while populating the enclave's page table with entries.
The kernel module asks the SM to create one or several threads and gives them to the enclave. 
The threads are set up to start executing at the right location and with the right stack pointer.
\\

\noindent\textbf{Timer Limit}
We also added a new value \texttt{timer\_limit} to each thread.
This value is used to avoid timing attacks on enclaves, especially ones exploiting malicious inputs to leak enclave secrets through the enclave completion time.
\texttt{timer\_limit} will represent the maximum amount of time the enclave's thread is allowed to execute.
During enclave launch, the hardware timers are set so the enclave will be exited after the corresponding amount of time has elapsed.
This way, if an unrecoverable exception occurs, the enclave will not immediately exit, but will simply wait for its timer to expire.
This prevents the enclave from leaking the timing at which exceptions occurred.
\\

\noindent\textbf{Sealing the Enclave}
The enclave can now be sealed, its measurement is finalized and used to generate a key pair.
The SM commits to the enclave's measurement and public key using its own signing key.
The enclave is now ready and the module can send a final ecall to the SM to enter the enclave.
\\

\noindent\textbf{Enclave Execution}
The enclave now executes on its dedicated cores and can communicate with its parent-level user-application using shared memory.
When the enclave wants to exit, it does so by sending a ecall directly to the mini-SM.
Upon return, the kernel module asks the SM to delete the threads and the enclave's metadata; it then reclaims the memory regions and gives them back to the OS.
Finally, the kernel module returns and forwards the enclave's return value to the user-level application.

\subsection{Enclave's Attestation Mechanism}
\label{app:attestation}
Remote attestation makes it possible for a remote user to receive and check a proof that convinces them of the integrity of an enclave's initial state, given trust in the Secure Bootloader and the SM's code.
Our remote attestation protocol is based on \cite{lebedev2018secure}.
The attestation proof is composed of two certificates, one ensuring the integrity of the SM and one ensuring the integrity of the secure enclave.
\\

\noindent\textbf{SM Attestation}
First, during secure boot, our (trusted) Secure Bootloader will measure the SM binary using a cryptographic hash function (SHA3) along with the device's secret, to derive the SM signing keys.
It will then generate a certificate using a public-key signature algorithm (Ed25519) with a key pair generated from a root device secret (generated on-chip, for example, using a PUF \cite{herder2014physical}) by signing the SM's measurement and its public key.
The SM certificate, the SM keys and the Bootloader public key are copied to SM memory by the Bootloader.
The SM certificate and the Bootloader public key are made available to the OS and to enclaves through the SM API.
The bootloader is then responsible for setting up the SM stack correctly and entering the SM at the right address.
\\

\noindent\textbf{Enclave Attestation}
During the enclave's setup, the SM will create the enclave's unique measurement by hashing (with SHA3) several arguments used to set up the enclave initial state.
In particular, it will include the enclave's trusted virtual address range, the enclave's binary along with its virtual memory mapping, the mini-SM binary and for each thread, the entry \texttt{PC} and \texttt{SP}, the value of \texttt{timer\_limit}, and a representation of which exceptions and interrupts are enabled and which ones are directly delegated to the enclave.
Because our SM is trusted, and the Ed25519 algorithms are implemented in constant time and protected against speculative and non-speculative side-channel attacks (speculation is disabled inside of our SM), we modify the original Sanctum scheme \cite{costan2016sanctum} and do {\em not} rely on an attestation enclave to sign the enclave's measurement and generate the enclave's certificate.

Instead, after the enclave's measurement has been completed, the SM will use it, along with a hash of its secret key, to derive a Ed25519 key pair for the enclave following a similar process to the Bootloader (see Figure \ref{fig:keyderiv}).
It will then generate an attestation certificate by signing the enclave's measurement and the enclave public key.
The enclave measurement, certificate and keys are copied to the enclave's metadata region, which is private to the SM.
\\

\noindent\textbf{Extending the SM Attestation API}
We introduce two new SM API calls.
One so an enclave can retrieve its own keys and attestation certificate (no one else can retrieve its private key).
And a second so anyone can retrieve any enclave attestation certificate (which contains its public key).
The enclave key pair can then be used to sign messages or to perform key exchanges and create secure communication channels \cite{bernstein2012high}.

This solution has several advantages compared to the Sanctum attestation enclave scheme.
The enclave does not need to rely on another enclave to perform its attestation.
This avoids several communication, concurrency and availability issues, for instance, how an enclave finds or wakes up the attestation enclave.
This also means we don't need to share the attestation enclave with other potentially untrusted enclaves.
Additionally, the key-derivation scheme is now deterministic.
This means that the same enclave on the same machine will be endorsed with the same keys given an identical SM and a fixed device secret.
This makes it possible for a remote user to process shared secrets with the enclave and prepare and encrypt data for the enclave without the enclave actively running.
This also makes it possible for an enclave to derive deterministic keys for long-term secure data storage that can be recovered even after being offline.
\\

\begin{figure}[t]
\centering
\includegraphics[width=\columnwidth]{figures/key_derivation.png}
\caption{Key Derivation and Endorsement}
\label{fig:keyderiv}
\end{figure}

\noindent\textbf{Position-Independent Enclave Measurement}
For our attestation mechanism to be practical and our enclave's key derivation to be deterministic, it is important for an enclave's measurement to be independent of the enclave's location in physical memory.
It would be impractical for a remote verifier to keep a copy of a hash for each possible enclave's placement in physical memory (given that an enclave could own any subset of the 64 physical memory regions, which would be at worst $2^{64}$ elements).
We need to tweak Sanctum's mechanism to properly enforce that property.
This property is enforced two ways.
\\
\noindent\textbf{Position-Independent Mini-SM}
First, for the mini-SM, its code must be completely position-independent. 
This is achieved at compilation time by using specific flags. 
Additionally, the mini-SM is only accessed through traps.
The entry into the trap handler and the mini-SM stack must be fixed relative to the mini-SM binary.
Because these values are physical addresses and cannot be added to the enclave's measurement, the SM is trusted to set up the trap handler and the stack correctly.
\\

\noindent\textbf{Mapping Virtual Addresses to Bytes}
Second, for the enclave's binary, we can rely on virtual memory to enforce that the enclave's code will always use the same addresses independently of its position in physical memory.
The only challenge left is how to provide a measurement that reflects the virtual address mapping without including any physical addresses. Instead of measuring the two mappings virtual-address-to-physical-address and physical-address-to-byte-content (for instance, by measuring the page table and measuring the physical memory content and its addresses), we will directly enforce and measure the mapping from virtual-addresses-to-byte-content.
When loading a page from the enclave's binary, the SM will be given, by the OS, the physical address at which the enclave page should be copied, the virtual address at which the page should be mapped and a pointer to the enclave's binary page in OS memory. 
It will then 1) copy the page to the destination physical address, 2) set up the enclave's page table to enforce the right mapping, and 3) only add to the enclave's measurement the hash of the page content, the hash of the virtual address it was mapped to, and the hash of the page table entry flags used.
As a result, it is enough for a remote verifier that trusts the SM to 1) check the SM certificate and 2) check the enclave's measurement to correctly identify the enclave.
\\

\noindent\textbf{Extension to Trusted VMs}
Note that for an enclave that manages its own virtual memory, we would need to play other tricks to make its measurement position independent.
Indeed, the enclave would have somehow to learn about its position in physical memory, by hardcoding it (and hence fixing it) it or by reverse-engineering the page table provided by the SM.
Hence, porting a complete operating system inside an enclave would require modifications to achieve a realistic end-to-end attestation solution.

\subsection{Building an Enclave Application}

\noindent\textbf{Enclave Runtime}
The enclave binary or enclave application is the code that will run inside the enclave.
In should really be seen as a binary as, at first, the execution environment made available to the enclave is close to bare metal.
Indeed, to enforce the integrity and privacy of the enclave, its private memory and its execution cores are architecturally and microarchitecturally isolated from the OS.
This means, for instance, that the enclave cannot directly interact with the OS through the usual syscalls.
As a result, an enclave binary needs to include its own libraries.
Recall also that the OS does not have access to the enclave's private memory.
That means the enclave doesn't have access to the OS resource management and needs to manage resources (e.g., memory, threads, etc.) on its own.

However, the enclave has access to its mini-SM through the \verb|ecall| interface and can perform operations such as freeing a memory region, setting up timer interrupts, or exiting.
The enclave also has access to shared memory: it can read and write the data that is made available to it by the OS.
In our implementation, the enclave can access the address space of the user-level application that launched it.
This makes our enclaves expressive and allows them to run rich applications such as cryptographic libraries, language run-times, or ML inference.
Inputs can be shipped directly with the enclave binary or communicated to the enclave using shared memory.
To protect the integrity and privacy of the input data sent to an enclave, a remote user can establish a shared secret with the enclave using the enclave's public key and encrypt the data before sending them to the enclave (see Section \ref{app:attestation}).


\subsection{Debugging Infrastructure}
\noindent\textbf{Debugging an Application}
For early software development and debugging, we developed a QEMU target that matches our processor.
A programmer can easily run their applications on several cores on top of QEMU without requiring to run the enclave from within Linux.
From there, the application can be easily debugged using GDB.
We also make it possible to compile the mini-SM with specific debugging flags to allow an enclave to (insecurely) access the console and to have the SM printing extra debugging messages, such as information regarding enclave exceptions.
Once Linux has been booted, applications can be debugged on FPGA using standard Linux debugging tools.
\\

\noindent\textbf{Debugging the Enclave Platform}
Our QEMU simulator can also be used to debug other elements of the software stack including the Secure Bootloader, the Security Monitor, Linux itself, but also the SM Kernel Module, enclave binaries, and user-level applications.
In addition, the hardware can be compiled to run in Verilator, a cycle-accurate simulator.
This can be especially handy for debugging new hardware features or bugs that arise only on the FPGA.
We also developed a series of low-level hardware tests to check the memory protection primitives such as the dual-page-table mechanism.
To conclude, the debugging infrastructure makes it possible to efficiently implement new mechanisms whether they are implemented in hardware, software or a combination of both.